\section{From Award-Layer Triage to Bid-Layer Forensics}
\label{sec:screening_forensics}

The preceding sections validate the award-layer ranking on the
margin for which it is designed. Section~\ref{sec:validation}
shows that persistent zero-win participation orders
adjudication-anchored loser-side exposure; Section~\ref{sec:cobidder_type}
shows that the exposed firms have an economically distinctive
profile rather than merely high procurement volume. The remaining
question is how an enforcement agency should use such a signal.
The answer is not to stop at the award layer, and it is not to
replace bid-layer forensics. A thin administrative record can
rank priorities, but it cannot by itself evaluate the bid conduct
that richer cartel evidence requires.

% COMPRESS EDIT (length): collapsed the three framing paragraphs and the
% bid-layer benchmark subsection (5 floats, feature-by-feature/fold-by-fold
% narration) into two compact paragraphs + Table tab:bid_layer_performance and
% (Subprompt 10) the cost-recall frontier Table tab:cost_recall_frontier;
% full bid-benchmark detail lives in Appendix~I.
This section turns from validation to sequencing. The screen is
useful because it sits before costly proof-producing work: it
uses records the agency observes broadly---who participated, who
won, and who kept appearing without winning---whereas bid-layer
forensics use submitted offers, within-tender bid moments, and
distance from the winning bid, which are more probative but more
expensive to recover, clean, and inspect at scale. The issue is
not whether a cheap signal should replace a costly one, but
whether it can discipline where the costly one is opened first.
We evaluate that architecture in three steps: we place the
award-layer ranking and an Imhof--Wallimann-style bid-distribution
benchmark on the same adjudication-anchored target, so performance
can be read against information cost; we ask whether the two
layers are complements; and we implement the operational
gatekeeper, in which award records create a survivor pool and bid
forensics rank firms inside it. The object is a cost-recall
frontier for forensic attention, not a horse race under full
observability. This is the paper's central enforcement
implication: administrative records rank where proof-producing
forensics should begin, while bid-level evidence remains the stage
at which conduct is evaluated.

\subsection{Bid-Layer Forensics as a Full-Observability Benchmark}
\label{sec:bid_layer_benchmark}

The two layers sit at different points in the evidentiary
pipeline. The award layer uses records the agency already holds;
the bid layer uses the submitted offers, which in the Brazilian
setting are not in the routine award file and must be recovered
from the underlying \texttt{LANCES} bid histories, harmonized, and
reduced to within-tender moments before any firm can be scored.
The bid-layer benchmark is therefore the forensic stage, not a
free comparator: it becomes available only after the agency has
paid the cost of opening the bid file. The right question is how
much ranking information is already present before that file is
opened.

The benchmark is a fair, transparent forensic comparator rather
than a black box. The learner is a \valBidLearner{} (ranger, 500
trees, probability mode, no hyperparameter search), trained on
seven within-tender Imhof--Wallimann moments---the coefficient of
variation and its within-firm dispersion, skewness, kurtosis,
spread, the log min--max ratio, and the second-lowest-bid
gap---each averaged across the tenders a firm contests
\citep{abrantesmetz2006a,imhof2018screening,imhof2019detecting,huber2019machine,wallimann2023machine}.
The candidate pool is the always-losers for which these moments
are computable, the target is the same adjudication-anchored
cobidder exposure used throughout Section~\ref{sec:validation},
and all $\valBidPosRetained$ positives are retained; the small
attrition is mechanical and concentrated among non-positives.
Timing and bid-revision features are excluded because the records
carry no usable per-bid timestamps. The candidate-support audit,
the model-design audit, and the leakage-boundary table are
reported in full in Appendix~\ref{app:bid_benchmark_submission}.

\paragraph{Performance.}
Table~\ref{tab:bid_layer_performance} reports precision, recall,
and PR-AUC first and ROC second, since the positive base rate is
near one percent and ranking quality is what an agency cares
about. Two facts stand out. First, in the pooled
full-observability design the cheap award flag is \emph{comparable}
to the costly bid benchmark: the binary frequent-loser flag
reaches ROC $\valImhofFLBin$ and the bid-layer \valBidLearner{}
reaches ROC $\valImhofFull$, with overlapping confidence
intervals---comparable discrimination at lower informational
cost, not a victory of one layer over the other. The combined
award-plus-bid model reaches ROC $\valImhofComboCont$ and
PR-AUC $0.275$, but it requires bid microdata for every firm and
is therefore a full-observability upper bound rather than a
deployable screen. Second, the pooled metrics rest on
\emph{random} cross-validation folds that are optimistic because
the cobidder positives cluster by CADE case. Under case-grouped
folds, where the positives from a held-out case never co-train,
the bid layer falls hardest: its PR-AUC drops from
$\valBidPooledImhofPRAUC$ to $\valBidLOCOimhofPRAUC$, a
$\valBidLOCOimhofPRAUCdrop$ reduction, and its ROC from $0.891$ to
$\valBidLOCOimhofAUC$. The combined model holds up (ROC
$\valBidLOCOcomboAUC$) only because the label-independent award
score continues to rank firms when the bid signal thins.
Excluding the label-defining tenders lowers the bid ROC only
modestly (to $\valBidExclLabelAUC$, $-0.017$), so the fragility is
case clustering, not feature--target overlap.

\paragraph{Complementarity.}
The two layers carry partially different ranking information
rather than duplicating one another: the Spearman rank correlation
between the award and bid scores is only $0.42$, and adding the
award score to the bid benchmark improves PR-AUC by
$+0.151$ under pooled folds (DeLong $p<0.001$). We report the
honest bound on that increment: under case-grouped folds it more
than halves to $+0.064$ PR-AUC (still $p<0.001$), and what
survives is mostly the award score \emph{rescuing} the bid layer
rather than the reverse. The combined model is therefore an upper
bound on what the two layers can jointly do under full
observability, not an operational first-stage rule. Neither score
is proof of an agreement: the bid benchmark measures the shape of
a firm's bid distribution, and false positives are non-labeled
firms, not ``bad firms.'' A strict-timing version of the bid
benchmark is \emph{blocked}, because within-tender dispersion
moments cannot be cleanly restricted to a pre-target window; the
award score, by contrast, is label-independent and admits a clean
temporal holdout. The full complementarity decomposition,
rank-overlap diagram, and leakage boundary appear in
Appendix~\ref{app:bid_benchmark_submission}. What the benchmark
establishes is the discriminating ceiling under full
observability, and how much of it the cheap award layer already
reaches at far lower informational cost.

\begin{table}[htbp]
\centering
\caption{Bid-Layer Benchmark and Award-Layer Screen: Ranking Performance}
\label{tab:bid_layer_performance}
\begin{adjustbox}{max width=\textwidth,center}
\begin{threeparttable}
\scriptsize
\begin{tabular}{llrrrrr}
\toprule
Design & Model & PR-AUC & Prec@500 & Rec@500 & Prec@1000 & ROC \\
\midrule
\multirow{4}{*}{Pooled (random folds)}
 & Award continuous & $0.126$ & $0.130$ & $0.342$ & $0.095$ & $0.938$ \\
 & Award FL14 & $0.085$ & $0.088$ & $0.232$ & $0.089$ & $\valImhofFLBin$ \\
 & Bid-layer \valBidLearner{} & $\valBidPooledImhofPRAUC$ & $0.136$ & $0.358$ & $0.100$ & $\valImhofFull$ \\
 & Combined award $+$ bid & $0.275$ & $0.198$ & $0.521$ & $0.149$ & $\valImhofComboCont$ \\
\midrule
\multirow{4}{*}{Case-grouped (held-out case)}
 & Award continuous & $0.126$ & $0.130$ & $0.342$ & $0.095$ & $0.938$ \\
 & Award FL14 & $0.085$ & $0.088$ & $0.232$ & $0.089$ & $\valImhofFLBin$ \\
 & Bid-layer \valBidLearner{} & $\valBidLOCOimhofPRAUC$ & $0.054$ & $0.142$ & $0.059$ & $\valBidLOCOimhofAUC$ \\
 & Combined award $+$ bid & $\valBidLOCOcomboPRAUC$ & $0.122$ & $0.321$ & $0.121$ & $\valBidLOCOcomboAUC$ \\
\midrule
Excl.\ label-defining tenders & Bid-layer \valBidLearner{} & $0.075$ & $0.092$ & $0.242$ & $0.084$ & $\valBidExclLabelAUC$ \\
\bottomrule
\end{tabular}
\begin{tablenotes}
\scriptsize
\item \textit{Notes:} $N=16{,}731$ always-losers with computable
bid moments; $190$ in-pool cobidder positives. PR-AUC is reported
first because the positive base rate is near one percent. The
\textbf{pooled} block is a full-observability diagnostic with
random-fold optimism; the \textbf{case-grouped} block is the
honest robustness test (positives clustered by case never
co-train) and shows the bid layer falling hardest (PR-AUC
$\valBidPooledImhofPRAUC\!\to\!\valBidLOCOimhofPRAUC$,
$\valBidLOCOimhofPRAUCdrop$). The award scores are
label-independent and so do not change across designs. Cobidders
are adjudication-anchored exposure, not cartel membership; false
positives are non-labeled firms, not ``bad firms.'' The combined
model is a full-observability upper bound, not a first-stage
screen. Candidate support, model audit, complementarity, and the
leakage boundary are in
Appendix~\ref{app:bid_benchmark_submission}.
\end{tablenotes}
\end{threeparttable}
\end{adjustbox}
\end{table}

% CO-AUTHOR EDIT (Subprompt 10): replaced the single-point gatekeeper subsection
% (former tab:gatekeeper_submission, \valArch* macros) with a cost-recall frontier
% subsection built on Table~\ref{tab:cost_recall_frontier} and Fig~\ref{fig:cost_recall_frontier}.
% The contribution is the frontier, not a calibrated K1; the firm-vs-bid-row divergence
% is reported as the honest core. The old "Timing the Gatekeeper" subsection is folded
% into one compressed paragraph below.
\subsection{Sequential Gatekeeping and the Cost-Recall Frontier}
\label{sec:sequential_gatekeeper}

Similar discrimination at lower cost does not by itself establish
a useful two-stage architecture. The two layers sit at different
points in the evidentiary pipeline, and that ordering is the
operational lever. The award layer is routine and cheap: it uses
records the agency already holds. The bid layer is richer and more
probative but costlier: in this setting the submitted offers are
not in the award file and must be recovered from the underlying
\texttt{LANCES} histories, harmonized, and reduced to within-tender
moments before any firm can be scored. A model that gives the
classifier both layers for every firm---the \emph{joint} model---is
an upper bound on what the two layers can do under full
observability, but it has already paid the full cost of opening the
bid file. The implementable object is therefore sequential: rank
where the costly layer should be opened first.

The sequential rule imposes that constraint directly. Firms are
ranked using only the award layer; the top $K_1$ firms form a
survivor pool $S_{K_1}$; bid microdata are recovered \emph{only}
for the survivors in $S_{K_1}$; and the bid-distribution benchmark
then reranks firms within that smaller pool by the combined
award-plus-bid score, producing the final top-$k$ forensic queue.
Full observability is the upper-bound benchmark, not the deployed
policy; award-only is the zero-bid-cost baseline that opens no bid
records at all; the sequential rule is the operating architecture
in between.

\paragraph{Cost is not only firms.}
The first-order honest point is that the firm count is the wrong
denominator for forensic burden. Recovering the within-tender bid
moments requires the full tender-item \texttt{LANCES} record for
every tender a survivor contested, so the relevant cost is measured
in opened tender-items and opened bid rows, not in firms.
Table~\ref{tab:cost_recall_frontier} reports all three. At the
$K_1=2{,}000$ operating point the rule opens only $12\%$ of firms
(a firm reduction of $\valCostFirmRedTwoK$ relative to full
observability), but it
recovers $\valCostTenderRedTwoK$ fewer tender-items and only
$\valCostBidRowRedTwoK$ fewer bid rows. Survivors are the
highest-participation always-losers, so opening $12\%$ of firms
still opens roughly $67\%$ of the bid rows. The frontier reports
firms, tender-items, and bid rows precisely so this divergence is
visible: the honest processing-burden cut is about a third at the
bid-row level, not the firm-level figure. We report the
firm-vs-bid-row divergence ($\valCostFirmRedTwoK$ versus
$\valCostBidRowRedTwoK$) as the core caveat, not a footnote.

\paragraph{The frontier, not a cutoff.}
Varying $K_1$ traces a cost-recall frontier rather than identifying
a single best rule. At $k=500$, sequential award$\rightarrow$bid
recall is non-monotone in $K_1$:
$0.34$ at $K_1=500$, $0.44$ at $1{,}000$,
$\valCostRecallSeqTwoKkFiveH$ at $2{,}000$,
$\valCostRecallSeqThreeKkFiveH$ at $3{,}000$, and $0.42$ at $5{,}000$,
peaking near $K_1=3{,}000$ (Fig.~\ref{fig:cost_recall_frontier}).
At $K_1=3{,}000$ the sequential rule \emph{matches} joint
full-observability recall ($\valCostRecallSeqThreeKkFiveH$ versus
$\valCostRecallJointkFiveH$ at $k=500$) at roughly a $24\%$ bid-row
reduction, but it never exceeds the joint upper bound
($\valCostRecallJointkFiveH$ at $k=500$,
$\valCostRecallJointkOneK$ at $k=1{,}000$). $K_1=2{,}000$ is
\emph{one operating point, not a calibrated optimum}: agencies with
larger forensic capacity sit higher on the frontier and agencies
with less capacity sit lower. The contribution is the frontier
itself---the menu of recall-versus-scope trade-offs---not a
universal cutoff.

\paragraph{What the bid rerank does, and does not, add.}
The bid layer adds little operationally at a fixed survivor pool.
At $K_1=2{,}000$ the award-defined survivor pool \emph{already}
contains $\valCostSurvRecallTwoK$ of the positives the final queue
recovers; the subsequent bid rerank only reorders firms inside that
pool and does not surface new positives. The architecture's value
is therefore in cost---preserving much of the target capture while
reducing the scope of proof-producing forensics---rather than in
the bid layer locating targets the award layer missed. This is the
opposite of a claim that the bid stage drives recovery: it is the
award gate that determines which positives are reachable, and the
rerank that orders them.

\paragraph{False positives and missed positives.}
The frontier is explicit about both error margins. At
$K_1=2{,}000$, $k=500$, the rule opens $\valCostFPrecoveryTwoK$
non-cobidder firms for bid recovery, loses
$\valCostPosLostGateTwoK$ positives at the award gate (they never
enter the survivor pool), and ends with $409$ non-labeled firms in
the final top-$500$ queue; $\valCostMissedTwoKkFiveH$ of the
$\valCostNpos$ positives are missed overall, and precision is
modest at about $0.18$. We do not read modest precision as failure:
in forensic triage, high recall that routes a manageable queue to
the next, costlier evidentiary stage is the operative property, and
the $409$ non-labeled firms in the queue are unverified, not
exonerated.

\paragraph{Baselines.}
The frontier is bracketed by three reference points. Award-only is
the \emph{zero-bid-cost} baseline (recall
$\valCostRecallAwardkFiveH$ at $k=500$), achievable with no bid
recovery whatever. Bid-only and joint scoring are full-observability
benchmarks that require opening the entire pool's bid records; joint
is the upper bound the sequential rule approaches but never
exceeds. A random survivor pool of the same size is far below the
award-ranked one (recall $\valCostRandomRecallTwoK$ at
$K_1=2{,}000$, three to twelve times worse across operating points),
confirming that it is the award ranking, not the mere act of opening
$K_1$ firms, that carries the signal.

The institutional reading follows directly. The rule is not
designed to beat full information; it preserves much of the target
capture while reducing the proof-producing forensic scope---and
even that reduction is honest only at the firm level and modest at
the bid-row level. The bid rerank's marginal contribution at a
fixed survivor pool is reordering, not new recall. A screen of this
kind does not decide liability; it decides, under a stated capacity
constraint, which files deserve the next and more expensive form of
evidentiary attention.

The frontier is also case-fragile. Leaving out the single largest
CADE case---which supplies $55\%$ of the positives---drops
sequential recall at $K_1=2{,}000$ from $\valCostRecallSeqTwoKkFiveH$
to $\valCostLOCOdrop$ at $k=500$. The signal is real but
concentrated; a different case mix would move the frontier, and the
operating points should be read as setting-specific rather than
portable magnitudes.

\begin{table}[htbp]
\centering
\caption{Cost-Recall Frontier for Sequential Award$\rightarrow$Bid Gatekeeping}
\label{tab:cost_recall_frontier}
\begin{adjustbox}{max width=\textwidth,center}
\begin{threeparttable}
\scriptsize
\begin{tabular}{lrrrrrrrrrrr}
\toprule
Rule & $K_1$ & Firms & Tender & Bid & TP & FP & Recall & Prec. & $\Delta$Firm & $\Delta$Bid & RLoss \\
 & & open & items & rows & & & & & red. & red. & vs joint \\
\midrule
\multicolumn{12}{l}{\textit{Panel A: forensic queue $k=500$}} \\
Award-only (no bid) & --- & $500$ & $41{,}704$ & $1{,}356{,}346$ & $65$ & $435$ & $\valCostRecallAwardkFiveH$ & $0.13$ & $0.97$ & $0.60$ & $0.18$ \\
Bid-only (full obs) & --- & $\valCostPoolN$ & $116{,}112$ & $3{,}393{,}915$ & $66$ & $434$ & $0.35$ & $0.13$ & --- & --- & $0.17$ \\
Joint (full obs, U.B.) & --- & $\valCostPoolN$ & $116{,}112$ & $3{,}393{,}915$ & $99$ & $401$ & $\valCostRecallJointkFiveH$ & $0.20$ & --- & --- & $0.00$ \\
Sequential & $500$ & $500$ & $41{,}704$ & $1{,}356{,}346$ & $65$ & $435$ & $0.34$ & $0.13$ & $0.97$ & $0.60$ & $0.18$ \\
Sequential & $1{,}000$ & $1{,}000$ & $58{,}408$ & $1{,}783{,}750$ & $83$ & $417$ & $0.44$ & $0.17$ & $0.94$ & $0.47$ & $0.08$ \\
Sequential & $2{,}000$ & $2{,}000$ & $77{,}097$ & $2{,}283{,}924$ & $91$ & $409$ & $\valCostRecallSeqTwoKkFiveH$ & $0.18$ & $\valCostFirmRedTwoK$ & $\valCostBidRowRedTwoK$ & $0.04$ \\
Sequential & $3{,}000$ & $3{,}000$ & $87{,}714$ & $2{,}575{,}274$ & $99$ & $401$ & $\valCostRecallSeqThreeKkFiveH$ & $0.20$ & $0.82$ & $0.24$ & $0.00$ \\
Sequential & $5{,}000$ & $5{,}000$ & $99{,}977$ & $2{,}921{,}977$ & $79$ & $421$ & $0.42$ & $0.16$ & $0.70$ & $0.14$ & $0.11$ \\
\midrule
\multicolumn{12}{l}{\textit{Panel B: forensic queue $k=1{,}000$}} \\
Award-only (no bid) & --- & $1{,}000$ & $58{,}408$ & $1{,}783{,}750$ & $95$ & $905$ & $0.50$ & $0.10$ & $0.94$ & $0.47$ & $0.28$ \\
Bid-only (full obs) & --- & $\valCostPoolN$ & $116{,}112$ & $3{,}393{,}915$ & $99$ & $901$ & $0.52$ & $0.10$ & --- & --- & $0.26$ \\
Joint (full obs, U.B.) & --- & $\valCostPoolN$ & $116{,}112$ & $3{,}393{,}915$ & $149$ & $851$ & $\valCostRecallJointkOneK$ & $0.15$ & --- & --- & $0.00$ \\
Sequential & $1{,}000$ & $1{,}000$ & $58{,}408$ & $1{,}783{,}750$ & $95$ & $905$ & $0.50$ & $0.10$ & $0.94$ & $0.47$ & $0.28$ \\
Sequential & $2{,}000$ & $2{,}000$ & $77{,}097$ & $2{,}283{,}924$ & $127$ & $873$ & $\valCostRecallSeqTwoKkOneK$ & $0.13$ & $\valCostFirmRedTwoK$ & $\valCostBidRowRedTwoK$ & $0.12$ \\
Sequential & $3{,}000$ & $3{,}000$ & $87{,}714$ & $2{,}575{,}274$ & $140$ & $860$ & $0.74$ & $0.14$ & $0.82$ & $0.24$ & $0.05$ \\
Sequential & $5{,}000$ & $5{,}000$ & $99{,}977$ & $2{,}921{,}977$ & $113$ & $887$ & $0.59$ & $0.11$ & $0.70$ & $0.14$ & $0.19$ \\
\bottomrule
\end{tabular}
\begin{tablenotes}
\scriptsize
\item \textit{Notes:} Evaluation pool: $\valCostPoolN$ always-loser
firms with computable complete-Imhof bid-distribution features;
positives are $\valCostNpos$ CADE-cobidders. The target is
\emph{adjudication-anchored exposure}, not cartel membership.
``Firms open,'' ``Tender items,'' and ``Bid rows'' report the
forensic footprint that must be recovered before the rule's final
queue can be scored; full observability opens all $\valCostPoolN$
firms / $116{,}112$ tender-items / $3{,}393{,}915$ bid rows.
\textbf{The firm count is not the only denominator}: opened
tender-items and bid rows better approximate processing burden, and
the firm reduction ($\Delta$Firm red.) substantially overstates the
bid-row reduction ($\Delta$Bid red.) because survivors are the
highest-participation firms. \textbf{Award-only opens no bid records}
as a pure award screen; the bid-footprint columns for the award-only
rows are the \emph{hypothetical} scope if it were verified at the
bid layer. Bid-only and joint require full bid observability for the
entire pool; the sequential rule recovers bid microdata only for the
$K_1$ survivors. Cost reductions are computed relative to full
observability. ``RLoss vs joint'' is recall lost relative to the
joint upper bound at the same $k$. False positives (FP) are
non-labeled firms, \emph{not} exonerated. ROC is not the decision
metric here; recall against the adjudication-anchored target at a
fixed forensic queue size is.
\end{tablenotes}
\end{threeparttable}
\end{adjustbox}
\end{table}

\begin{figure}[htbp]
\centering
\includegraphics[width=0.82\textwidth]{./output/figures/fig_3_cost_recall_frontier}
\caption{Cost-recall frontier for sequential award$\rightarrow$bid
gatekeeping. Points trace operational trade-offs in forensic
attention, \emph{not} legal proof. The horizontal axis is the
recovered bid-record footprint (bid-row reduction relative to full
observability) and the vertical axis is recall against
adjudication-anchored cobidder exposure at a fixed forensic queue
size. Full-observability benchmarks (bid-only, joint) require
whole-pool bid recovery and are plotted as horizontal reference
lines; the sequential points vary $K_1\in\{500,1{,}000,2{,}000,
3{,}000,5{,}000\}$ and never exceed the joint upper bound;
award-only is the zero-bid-cost baseline.}
\label{fig:cost_recall_frontier}
\par\noindent\textit{Alt text:} A scatter-and-line plot. The
horizontal axis is the reduction in recovered bid records relative
to opening every firm's bid history; the vertical axis is recall of
adjudication-anchored cobidder positives at a fixed forensic queue
size. A curve of sequential operating points rises and then falls as
$K_1$ grows from 500 to 5{,}000, peaking near $K_1=3{,}000$ where it
touches a horizontal dashed line marking the joint full-observability
upper bound; it never rises above that line. A second horizontal line
marks the lower bid-only full-observability benchmark, and a marker at
the far right shows the zero-bid-cost award-only baseline at lower
recall. The visual message is that opening fewer bid records costs
recall along a frontier, the sequential rule at best matches but never
exceeds full information, and the firm-level cost saving overstates the
bid-row cost saving.
\end{figure}

\paragraph{Timing discipline.}
The frontier above uses the full feature window, which is close to
the adjudication target. A stricter timing audit restricts the
award-layer score to a 2009--2016 window and evaluates against later
adjudication-anchored exposure; this is not a real-time field
deployment, since the labels remain adjudication-based, but it
disciplines the information used to rank firms before bid forensics
are allocated. Under that restriction the \emph{award-only} score
remains usable, with ROC falling to $0.73$ from the full-window
$0.94$. The \emph{sequential} rule, by contrast, cannot be cleanly
time-limited: the within-tender bid moments that drive the rerank
cannot be restricted to a pre-target window without contaminating
them with post-window tenders, so a strict-timing sequential frontier
is \emph{blocked}. We state this rather than fabricate a number. The
conclusion of Section~\ref{sec:screening_forensics} is therefore
operational and bounded: cheap administrative records can structure
the order in which expensive proof-producing evidence is assembled
and, on the award layer, survive a timing holdout; the richer bid
layer remains the place where conduct is evaluated, and its strict
timing properties cannot be certified here.
